\section{子任务 1：第 0 章 Introduction}

\subsection{核心知识点}

\begin{itemize}[leftmargin=2em]
  \item 给出全书三大主线：automata、computability、complexity。
  \item 统一后续讨论所需的数学语言：集合、序列、关系、图、字符串、语言、布尔逻辑。
  \item 强调定义、定理、证明的角色分工：定义负责划边界，定理负责陈述规律，证明负责提供绝对说服力。
  \item 介绍三种最基础的证明方式：构造法、反证法、归纳法。
\end{itemize}

\subsection{典型问题与证明套路}

\begin{enumerate}[leftmargin=2em]
  \item \textbf{要证明某对象存在}：优先用构造法，直接把对象写出来。
  \item \textbf{要证明某陈述不可能}：优先尝试反证法，假设存在再推出矛盾。
  \item \textbf{要证明一个与长度、步数、层数相关的性质}：优先尝试归纳。
  \item \textbf{要判断一个定义是否好用}：看它能否精确区分“属于”和“不属于”的对象。
\end{enumerate}

\subsection{证明背后的哲学}

第 0 章真正要训练的不是技巧，而是口味。Sipser 一开始就把“好证明”和“坏证明”的区别讲清楚：数学里不接受“差不多对”，也不接受“看起来像”。一条定义必须能承受后续几十页推理的压力，因此精确性不是形式主义，而是整个理论结构的承重墙。

\subsection{本章精华}

\begin{itemize}[leftmargin=2em]
  \item 计算理论不是从机器开始，而是从\textbf{如何严密地谈论机器}开始。
  \item 后面所有大定理，其实都建立在第 0 章这套定义与证明纪律之上。
\end{itemize}

\newpage
